\clearpage\FloatBarrier\section{Constructed graph schedules, support, and preconditioning}\label{supp5:section:001}
This appendix develops support bounds, graph schedules, block discovery and preconditioning certificates. Classical algorithmic inputs remain separate from the finite certificates proved here.\par

\subsection{Constructing the optimal level schedule}\label{supp5:subsection:001}\textit{Source items 6.6 and 10.3; module GraphSchedules}\par
Let $V$ be finite and let $u\to v$ mean that task $v$ depends on task $u$. Assume there is no nonempty directed cycle. A topological order or valid level function is \emph{not} supplied as an additional hypothesis.

For each vertex, define its strict ancestor set
\[
\operatorname{Anc}(v)=\{u:\text{there is a positive-length path }u\leadsto v\}.
\]
For every edge $u\to v$, the proof establishes
\[
\operatorname{Anc}(u)\subsetneq\operatorname{Anc}(v).
\]
The inclusion follows by appending the edge. It is strict because $u$ is an ancestor of $v$ but cannot be its own strict ancestor in an acyclic graph. Cardinality therefore decreases when following predecessors. This constructs a well-founded relation on which the level recursion is defined:
\[
\boxed{\ell(v)=1+\max_{u\to v}\ell(u),\qquad \max\varnothing=0.}
\]

\textbf{Validity and attainment.} Lean proves $\ell(v)>0$ and $\ell(u)<\ell(v)$ on every edge. A path of $n$ edges from $a$ to $b$ satisfies
\[
\ell(a)+n\leq\ell(b).
\]
Conversely, predecessor-max attainment and induction produce an actual path with $\ell(v)-1$ edges ending at each vertex $v$. The result is an equality with a longest path, not merely an upper bound.

A valid schedule $s:V\to\mathbb N$ assigns positive stages and strictly increases on each edge. Induction proves the stronger pointwise statement
\[
\ell(v)\leq s(v)\quad\text{for every }v.
\]
Consequently the constructed schedule achieves the least possible makespan,
\[
\boxed{d=\max_{v\in V}\ell(v).}
\]
For a nonempty graph, $d$ equals the number of vertices on an attained longest path. Vertices of the same level have no dependency edge in either direction.

\textbf{Empty graph and interpretation.} When $V=\varnothing$, the checked depth is zero. The phrase ``one plus the longest edge-path length'' applies to nonempty graphs. These are ideal discrete scheduling statements. Weighted task costs, a finite processor budget, launch latency and communication are not part of this model. The definitions are mathematical constructions, not a certified hardware scheduler.



\subsection{Forward substitution with an actual stage evaluator}\label{supp5:subsection:002}\textit{Source item 6.5; module TriangularSolve}\par
Use the graph and levels constructed on the previous page. A vertex $v$ carries an arbitrary vector space $X_v$ over a field. The spaces may have different dimensions. Let
\[
D_v:X_v\overset{\sim}{\longrightarrow}X_v,
\qquad T_{vu}:X_u\longrightarrow X_v
\]
be an invertible diagonal map and linear transfer maps. The operator being solved is defined by its actual block action:
\[
(Ax)_v=D_vx_v+\sum_{u\to v}T_{vu}x_u.
\]
No solution, solve equation or precomputed topological certificate is assumed.

\textbf{Constructed solution.} Well-founded recursion defines
\[
\boxed{x_v=D_v^{-1}\!\left(b_v-\sum_{u\to v}T_{vu}x_u\right).}
\]
Applying $D_v$ proves the original equation at $v$. If another vector satisfies all equations, induction over predecessors makes each previously used coordinate identical, and invertibility of $D_v$ forces agreement at $v$. Thus the construction gives existence and uniqueness, and the block operator is bijective.

\textbf{Parallel stages.} A second, independent definition is an explicit staged evaluator. Set $x^{(0)}=0$ and
\[
x^{(k+1)}_v=
\begin{cases}
D_v^{-1}\!\left(b_v-\displaystyle\sum_{u\to v}T_{vu}x^{(k)}_u\right),
  &\ell(v)=k+1,\\[5pt]
x^{(k)}_v,&\ell(v)\ne k+1.
\end{cases}
\]
It reads the previous-stage vector, not partially updated coordinates. The invariant is
\[
\ell(v)\leq k\quad\Longrightarrow\quad x^{(k)}_v=x_v.
\]
Every predecessor of a level-$k+1$ vertex has level at most $k$, so the recurrence uses only coordinates already proved correct. At the canonical depth $d$ this yields
\[
\boxed{Ax^{(d)}=b.}
\]
The local independence theorem separately proves that changing coordinates outside the strict earlier levels cannot change a vertex's local solve.

\textbf{Scope.} Diagonal invertibility is expressed by genuine linear equivalences. The task graph may include declared transfers whose values happen to be zero; optimality refers to that declared dependency graph. Exact field arithmetic and exact local inverses are used. No floating-point stability or implementation-specific parallel speed is inferred.



\subsection{Support of inverses and Schur elimination}\label{supp5:subsection:003}\textit{Source items 6.7 and 6.8; modules InverseSupport and SchurSupport}\par
Let $(P_a)$ be a complete idempotent packet family in a finite-dimensional algebra. Define $a\leadsto b$ by an actual finite path through nonzero blocks $P_bAP_a$. \textbf{Zero-length paths are included}, so every packet reaches itself.

Fix the source packet $a$. In the regular representation, define the submodule
\[
W_a=\{x:P_bx=0\ \text{whenever }a\not\leadsto b\}.
\]
It contains $P_a$. Completeness of the projectors expands each component of $Ax$ into packet transfers. If a transfer into an unreachable packet were nonzero from a reachable packet, it would extend a path, a contradiction. Thus $AW_a\subseteq W_a$.

Suppose $BA=I$. Left multiplication by $A$ is injective, including on $W_a$. Finite-dimensionality makes that restricted map surjective. Given $x\in W_a$, choose $y\in W_a$ with $Ay=x$. Then $Bx=BAy=y\in W_a$. Applying this to $P_a$ proves
\[
\boxed{a\not\leadsto b\quad\Longrightarrow\quad P_bBP_a=0.}
\]
The matrix specialization takes $B$ to Mathlib's actual nonsingular matrix inverse, with an invertible determinant. No inverse-sparsity hypothesis is inserted.

\textbf{Hidden elimination.} For the actual Schur complement
\[
S=A_{RR}-A_{RH}A_{HH}^{-1}A_{HR},
\]
the hidden inverse is expanded through the hidden projectors. A nonzero retained correction block requires indices $j,k$ with nonzero entrance and exit blocks and a hidden path $j\leadsto k$. Therefore a nonzero block $S_{ba}$ requires either a direct retained block, or
\[
a\longrightarrow j\ \leadsto_{H}\ k\longrightarrow b.
\]
The entrance and exit matrices may be rectangular; retained and hidden coordinate sizes and packet systems are independently typed. The hidden matrix must be invertible. A zero-length hidden path permits entrance and exit through the same hidden packet.

\textbf{Inclusion is not equality.} Cancellation can remove a block permitted by a graph route. For example,
\[
A=\begin{pmatrix}1&1&1\\0&1&1\\0&0&1\end{pmatrix}
\quad\text{has}\quad (A^{-1})_{13}=0
\]
despite a direct nonzero $(1,3)$ block. The implementation proves support inclusion and the corresponding no-route zero certificate, not a converse.



\subsection{Exact banded addresses and modeled work}\label{supp5:subsection:004}\textit{Source item 4.1; module BandedCount}\par
For $p$ packets and half-bandwidth $w$, the actual address set is
\[
\mathcal B_{p,w}
=\{(i,j)\in\{0,\ldots,p-1\}^2:
 i\leq j+w\ \text{and}\ j\leq i+w\}.
\]
The predicate is the integer version of $|i-j|\leq w$. The count is first defined as a double sum of indicators, then identified with the cardinality of this subtype. It is not defined by the target closed form.

Under $0\leq w<p$, the formal proof establishes
\[
\boxed{N_{p,w}=|\mathcal B_{p,w}|=p(2w+1)-w(w+1).}
\]
The proof uses finite prefix counting, the newly exposed boundary row and column when $p$ increases, and the saturated initial band. An intermediate subtraction-free identity avoids relying on natural-number subtraction without the required inequality. The diagonal case $w=0$ gives exactly $p$ blocks, and the saturated case $w=p-1$ gives $p^2$.

\textbf{Stored scalars.} With equal block width $b$, the storage-address type is
\[
\mathcal S=\mathcal B_{p,w}\times\{0,\ldots,b-1\}^2.
\]
Its cardinality is the actual number of represented scalar entries:
\[
\boxed{|\mathcal S|=b^2N_{p,w}.}
\]

\textbf{Arithmetic events.} For $q$ right-hand sides, the declared event type adds a right-hand-side index and one of two operation roles, multiplication or accumulation:
\[
\mathcal O=\{0,1\}\times\{0,\ldots,q-1\}\times\mathcal S.
\]
Hence
\[
\boxed{|\mathcal O|=2qb^2N_{p,w}.}
\]
These identities include zero right-hand-side count. The abstract storage formula also admits zero block width. To use the manuscript's substitution $p=n/b$, one must separately require $b>0$ and $n=pb$.

\textbf{Boundary.} Two events per coefficient and right-hand side are a declared arithmetic convention. These counts do not identify the number of CUDA instructions, account for fused multiply-add lowering, price complex arithmetic in real operations, or include allocation, indexing, communication and launch overhead. They are also not a performance comparison with dense multiplication.



\subsection{Exact residual certificates for supplied packets}\label{supp5:subsection:005}\textit{Source item 9.2; module DiscoveryCertificate}\par
Let $(P_a)$ be a complete orthogonal projector system on a finite complex coordinate space. The projectors need not be coordinate masks. For an actual matrix $A$, define
\[
\Pi(A)=\sum_aP_aAP_a,
\qquad E=A-\Pi(A).
\]
Reconstruction gives the exact off-diagonal residual
\[
E=\sum_{b\ne a}P_bAP_a.
\]

\textbf{Hilbert--Schmidt orthogonality.} For distinct ordered packet pairs $(b,a)\ne(d,c)$, adjoint multiplication and projector orthogonality give
\[
\operatorname{tr}\bigl((P_bAP_a)^*P_dBP_c\bigr)=0.
\]
If $b\ne d$, the central projector product vanishes. Otherwise $a\ne c$, and cyclicity of trace brings $P_cP_a$ together. This is proved for arbitrary input matrices $A,B$.

Lean identifies the actual Frobenius norm with its trace expression,
\[
\norm{M}_{\rm F}^2=\operatorname{Re}\operatorname{tr}(M^*M),
\]
then expands the finite sum and removes each cross term. The result is
\[
\boxed{\norm{A-\Pi(A)}_{\rm F}^2
=\sum_{b\ne a}\norm{P_bAP_a}_{\rm F}^2.}
\]
No Pythagorean equality for the target residual is assumed.

\textbf{Operator norm.} In a separate scoped matrix norm, the triangle inequality gives
\[
\boxed{\norm{A-\Pi(A)}_2
\leq\sum_{b\ne a}\norm{P_bAP_a}_2.}
\]
Here $\norm{\cdot}_2$ is the Euclidean induced operator norm, not the Frobenius norm. The formal declarations use distinct norm instances, preventing the two certificates from being conflated. Self-adjointness is used for Frobenius orthogonality; the sum estimate itself requires only the complete packet decomposition.

\textbf{Discovery boundary.} This result certifies the residual of a \emph{supplied} decomposition. It does not construct useful projectors from an arbitrary matrix, prove that a searched decomposition is optimal, or guarantee that its application is profitable on a GPU. The simultaneous spectral-decomposition input remains separately recorded as classical.



\subsection{Block-Jacobi normalization and the CG factor}\label{supp5:subsection:006}\textit{Source items 7.4 and 7.6; module BlockJacobi}\par
Work on a nonempty finite complex coordinate space with Euclidean operator norms. Assume $D$ and $A=D+K$ are positive definite. Define
\[
S=(\sqrt D)^{-1},\qquad B=SAS,\qquad E=SKS.
\]
The square root is the actual positive matrix square root supplied by continuous functional calculus; the inverse is the actual matrix inverse. Positive definiteness makes $\sqrt D$ a unit. Its square identity and inverse identities prove
\[
SDS=I,\qquad B=I+E.
\]

\textbf{A direct inverse bound.} If $C=(I+E)^{-1}$ and $\norm E_2\leq\eta<1$, then $C=I-EC$. Consequently
\[
(1-\eta)\norm C_2\leq 1,
\qquad \norm{I+E}_2\leq 1+\eta.
\]
The constructed normalization is invertible because $S$ and $A$ are. Combining the estimates gives the actual condition number bound
\[
\boxed{\kappa_2(B):=\norm B_2\norm{B^{-1}}_2
\leq\frac{1+\eta}{1-\eta}.}
\]
This norm argument does not require a separately assumed Hermitian perturbation. Nonnegative $\eta$ follows from the norm bound. Nonemptiness is explicit because the condition-number and norm-of-identity convention uses a nonzero coordinate space.

\textbf{The preconditioning corollary.} Define
\[
q(\kappa)=\frac{\sqrt\kappa-1}{\sqrt\kappa+1}.
\]
Monotonicity of square root and positive denominators prove that $q$ is monotone. Lean also proves $\kappa_2(B)\geq1$, $q(\kappa)\geq0$ for $\kappa\geq1$, and $q(\kappa)<1$. Thus
\[
\boxed{q(\kappa_2(B))\leq q\!\left(\frac{1+\eta}{1-\eta}\right).}
\]
For $1\leq\kappa$ and $e_0\geq0$, an explicit classical premise $e_m\leq2q(\kappa)^m e_0$ and a certified upper bound $\kappa\leq\bar\kappa$ yield $e_m\leq2q(\bar\kappa)^m e_0$.

\textbf{What is not claimed.} The classical conjugate-gradient convergence theorem, item 7.5, is not newly formalized. The new corollary proves its factor comparison and bound substitution. It does not replace the original theorem's algorithm hypotheses with a new axiom or claim a floating-point CG convergence proof.



\subsection{Place in the finite development}\label{supp5:subsection:007}
The global small-core recompression theorem is proved in Appendix F. The classical conjugate-gradient algorithm theorem, floating-point product model and simultaneous spectral decomposition remain separately supplied inputs.

The historical compiler, regression and development records appear in the formal companion under ``Historical proof-development records.
